% =============================================================================
% Section 5: Termal Deneme ve Basarisizlik Analizi
% INNATE Bitirme Tezi -- Berke Tezgocen
% =============================================================================

\section{Termal Deneme ve Basarisizlik Analizi}
\label{sec:thermal-failure}

Bu bolumde, INNATE modelinin izothermal (saf momentum) basarisinin ardindan
gerceklestirilen mixed convection genisletme denemesi detayli olarak
incelenmektedir. Termal deneme, bilimsel arastirma surecinin kacinilmaz bir
parcasi olan ``basarisiz deney'' kategorisinde yer almakla birlikte, elde
edilen bulgular modelin yapisal sinirlarini aydinlatmasi bakimindan en az
basarili sonuclar kadar degerli bilgi icermektedir. Sirasiyla motivasyon,
mimari genisletme, egitim surecinde karsilasilan sorunlar, kok neden analizi,
karsilastirmali perspektif ve alinan dersler ele alinmaktadir.


% =============================================================================
\subsection{Motivasyon ve Hedef}
\label{sec:thermal-motivation}

INNATE'in izothermal modda $\text{Re} = 7{,}000$ ve $\text{Re} = 10{,}000$'de
basarili sonuclar uretmesinin ardindan, dogal bir sonraki adim olarak
\emph{mixed convection} --- yani momentum ve isi transferinin birlesik
simulasyonu --- hedeflenmistir. Mixed convection, muhendislik uygulamalarinin
buyuk cogunlugunda (bina havalandirmasi, elektronik sogutma, atmosferik sinir
tabakasi) gecerli olan fiziksel rejimdir ve salt izothermal sonuclar bu
uygulamalara dogrudan aktarilamaz.

Mixed convection fizigin matematiksel cercevesi, Navier-Stokes denklemlerine
Boussinesq yaklasimi ile eklenen kaldirim kuvveti (buoyancy) ve enerji denklemi
uzerinden tanimlanir. Boyutsuz formulasyonda momentum denklemi:

\begin{equation}
  \frac{\partial \mathbf{u}}{\partial t}
  + (\mathbf{u} \cdot \nabla)\mathbf{u}
  = -\nabla p
  + \frac{1}{\text{Re}}\nabla^2 \mathbf{u}
  + \text{Ri}\,\theta\,\hat{\mathbf{e}}_y
  + \mathbf{F}
  \label{eq:momentum-boussinesq}
\end{equation}

\noindent burada $\text{Ri} = \text{Ra}/(\text{Re}^2 \cdot \text{Pr})$
Richardson sayisi, $\theta = (T - T_\text{ref})/\Delta T$ boyutsuz sicaklik
perturbasyonu ve $\hat{\mathbf{e}}_y$ yercekimi yonundeki birim vektordur.
Enerji (sicaklik) denklemi ise:

\begin{equation}
  \frac{\partial \theta}{\partial t}
  + (\mathbf{u} \cdot \nabla)\theta
  = \frac{1}{\text{Re}\cdot\text{Pr}}\nabla^2 \theta
  + v \cdot \frac{1}{L_y}
  \label{eq:energy-boussinesq}
\end{equation}

\noindent burada son terim ($v/L_y$), periyodik kutudaki ortalama sicaklik
gradyaninin baglantisidir (duvar sinir kosulu yoksa ``mean gradient'' kaynak
terimi).

Hedef buyuklukler su sekilde tanimlanmistir:

\begin{itemize}
  \item \textbf{Nusselt sayisi} ($\text{Nu}$): Konvektif isi transferinin
        iletimsel transfere orani.
        \begin{equation}
          \text{Nu} = 1 + \frac{\langle v\theta \rangle}{\kappa \cdot (1/L_y)}
          \label{eq:nusselt-def}
        \end{equation}
  \item \textbf{Termal flux} ($\langle v\theta \rangle$): Dikey hiz-sicaklik
        korelasyonu, turbulansi isi tasiminin dogrudan olcusu.
  \item \textbf{Sicaklik RMS} ($\theta_\text{rms}$): Sicaklik alaninin
        degiskenlik siddetini olcer.
        \begin{equation}
          \theta_\text{rms} = \sqrt{\langle \theta^2 \rangle}
          \label{eq:theta-rms-def}
        \end{equation}
\end{itemize}

LES (Large Eddy Simulation) referans degerleri Tablo~\ref{tab:les-thermal-ref}
ile verilmektedir.

\begin{table}[H]
  \centering
  \caption{LES referans termal metrikleri ($\text{Ra} = 10^5$, $\text{Pr} = 0.71$,
           $96 \times 160 \times 64$ grid, 60{,}000 adim).}
  \label{tab:les-thermal-ref}
  \begin{tabular}{lcccc}
    \toprule
    $\text{Re}$ & $\text{Nu}$ & $\langle v\theta \rangle$ & $\theta_\text{rms}$ & TKE \\
    \midrule
    7{,}000  & 39.3  & $7.71 \times 10^{-4}$ & 0.0240 & 0.00848 \\
    10{,}000 & 70.5  & $9.78 \times 10^{-4}$ & 0.0268 & 0.00731 \\
    \bottomrule
  \end{tabular}
\end{table}


% =============================================================================
\subsection{Termal Mimari Genisletme}
\label{sec:thermal-architecture}

INNATE'in 20 katmanli fractional-step mimarisine termal operatorler eklenerek
toplam parametre sayisi 167'den 287'ye cikarilmistir. Eklenen termal
``noronlar'' (ogrenilebilir fizik operatorleri) ve parametre dagilimi
Tablo~\ref{tab:thermal-params} ile ozetlenmektedir.

\begin{table}[H]
  \centering
  \caption{INNATE termal parametreleri (120 adet, toplam 287'nin \%42'si).}
  \label{tab:thermal-params}
  \begin{tabular}{llcl}
    \toprule
    Operator & Parametre & Sayi & Aciklama \\
    \midrule
    ThermalAdvection3D & \texttt{thermal\_adv\_modulator} & $20$ & Sicaklik adveksiyonu modulasyonu \\
    ThermalDiffusion3D & \texttt{kappa\_scale\_x/y/z} & $60$ & Anizotropik termal difuzyon \\
    EddyViscosity3D & \texttt{cs\_thermal} & $20$ & Termal SGS katsayisi \\
    Buoyancy3D & \texttt{buoyancy\_strength} & $20$ & Katman bazli kaldirim kuvveti \\
    \midrule
    \multicolumn{2}{l}{Toplam termal} & $120$ & \\
    \multicolumn{2}{l}{Toplam momentum} & $167$ & \\
    \multicolumn{2}{l}{\textbf{Genel toplam}} & $\mathbf{287}$ & \\
    \bottomrule
  \end{tabular}
\end{table}

Her katmandaki zaman adimi algoritmasi asagidaki siralamayi izler:

\begin{enumerate}
  \item \textbf{Advection:} $\mathbf{u}^* = \mathbf{u}^n - \Delta t \,
        (\mathbf{u}^n \cdot \nabla)\mathbf{u}^n + \Delta t\,\mathbf{F}
        + \Delta t\,\text{Ri}\,\theta^n\hat{\mathbf{e}}_y$
  \item \textbf{Diffusion:} $\mathbf{u}^{**} = \mathbf{u}^* + \Delta t\,
        \nu_\text{eff}\nabla^2\mathbf{u}^*$
  \item \textbf{Projection:} $\mathbf{u}^{n+1} = \mathbf{u}^{**} -
        \nabla p^{n+1}$, \quad $\nabla \cdot \mathbf{u}^{n+1} = 0$
  \item \textbf{Termal advection:} $\theta^* = \theta^n - \Delta t\,
        (\mathbf{u}^{n+1} \cdot \nabla)\theta^n$
  \item \textbf{Termal diffusion:} $\theta^{n+1} = \theta^* + \Delta t\,
        \kappa_\text{eff}\nabla^2\theta^* + \Delta t\, v^{n+1}/L_y$
  \item \textbf{Gauge fix:} $\theta^{n+1} \leftarrow \theta^{n+1}
        - \langle\theta^{n+1}\rangle$ (ortalama kaldirma)
  \item \textbf{Elevator mode removal:} $k_y = 0$ modlarini sifirla
\end{enumerate}

Termal egitim icin dort yeni loss terimi tanimlanmistir:

\begin{align}
  \mathcal{L}_{v\theta} &= \left(
    \log(|\langle v\theta\rangle| + \epsilon)
    - \log(|\langle v\theta\rangle_\text{ref}| + \epsilon)
  \right)^2
  \label{eq:loss-vtheta} \\
  \mathcal{L}_\text{corr} &= \left(
    \frac{\langle v\theta\rangle}{\sqrt{\text{TKE}} \cdot \theta_\text{rms}}
    - \frac{\langle v\theta\rangle_\text{ref}}{\sqrt{\text{TKE}_\text{ref}} \cdot \theta_{\text{rms,ref}}}
  \right)^2
  \label{eq:loss-corr} \\
  \mathcal{L}_{\theta,\text{rms}} &= \left(
    \log(\theta_\text{rms} + \epsilon)
    - \log(\theta_{\text{rms,ref}} + \epsilon)
  \right)^2
  \label{eq:loss-theta-rms} \\
  \mathcal{L}_{\theta,\text{min}} &= \left[
    \max(0,\; 0.015 - \theta_\text{rms})
  \right]^2
  \label{eq:loss-theta-min}
\end{align}

\noindent Nusselt referans loss'u ($\mathcal{L}_\text{Nu,ref}$) tanimlanmis
ancak \emph{damping paradoksu} nedeniyle sifir agirlikla devre disi
birakilmistir (Bolum~\ref{sec:damping-paradox}).


% =============================================================================
\subsection{Egitim Surecinde Karsilasilan Sorunlar}
\label{sec:training-problems}

Termal genisletme sonrasi gerceklestirilen coklu egitim denemeleri, asagida
kronolojik siralamayla ozetlenen sorunlarla karsilasmistir.

\subsubsection{Gradient Baskinligi}

Multi-objective loss fonksiyonunda momentum terimleri ($\mathcal{L}_\text{TKE}$,
$\mathcal{L}_\text{slope}$, $\mathcal{L}_\text{shape}$) termal terimlere
($\mathcal{L}_{\theta,\text{rms}}$, $\mathcal{L}_{v\theta}$) kiyasla
cok daha guclu gradient uretmektedir. Deneysel olcum:

\begin{table}[H]
  \centering
  \caption{Loss terimlerinin gradient norm karsilastirmasi (tipik egitim).}
  \label{tab:grad-norms}
  \begin{tabular}{lcc}
    \toprule
    Loss terimi & $\|\nabla_\theta \mathcal{L}\|$ & Sinif \\
    \midrule
    $\mathcal{L}_\text{TKE,ref}$       & $\sim 1{,}153$ & Momentum \\
    $\mathcal{L}_\text{slope,ref}$     & $\sim 487$     & Momentum \\
    $\mathcal{L}_\text{shape}$          & $\sim 312$     & Momentum \\
    $\mathcal{L}_{\theta,\text{rms}}$   & $\sim 12$      & Termal \\
    $\mathcal{L}_{v\theta}$             & $\sim 8$       & Termal \\
    $\mathcal{L}_\text{corr}$           & $\sim 5$       & Termal \\
    $\mathcal{L}_\text{Nu,ref}$         & ---            & Devre disi \\
    \bottomrule
  \end{tabular}
\end{table}

Bu dengesizlik, multi-objective optimizasyonda bilinen \emph{gradient dominance}
sorununa karsilik gelmektedir. Toplam loss:

\begin{equation}
  \mathcal{L}_\text{total}
  = \sum_{i \in \mathcal{M}} w_i \mathcal{L}_i
  + \sum_{j \in \mathcal{T}} w_j \mathcal{L}_j
  \label{eq:total-loss}
\end{equation}

\noindent icin, momentum terimleri domine ettiginde optimizer efektif olarak
sadece $\nabla_\theta \sum_{i \in \mathcal{M}} w_i \mathcal{L}_i$ yonunde
hareket eder. Termal parametreler (ozellikle \texttt{cs\_thermal},
\texttt{kappa\_scale}) uzerindeki gradient sinyali gurultu seviyesine duser.

Bu sorunu hafifletmek icin \emph{gradient routing} stratejisi uygulanmistir:
parametreler momentum, termal ve ``bridge'' gruplarina ayrilmis, her grubun
sadece ilgili loss terimlerinden gradient almasi saglanmistir. Ancak bu
yaklasim momentum-termal coupling'in dogasi geregi tam olarak izole edilemez
(Bolum~\ref{sec:vtheta-deadlock}).


\subsubsection{$\theta_\text{rms}$ Cokusu (Isothermal Collapse)}
\label{sec:theta-collapse}

Egitim suresinin ilk 100--200 epoch'unda, sicaklik alani sistematik olarak
duzlesmekte ($\theta_\text{rms} \to 0$) ve model fiilen izothermal moda
geri donmektedir. Bu davranis ``isothermal collapse'' olarak
adlandirilmistir.

Mekanizma su sekilde gerceklesmektedir: Termal gradient zayiflayinca
buoyancy terimi $\text{Ri}\,\theta\hat{\mathbf{e}}_y \to 0$ olur, bu da
momentum denklemini saf izothermal hale indirir. Izothermal modda momentum
loss'lari ($\mathcal{L}_\text{TKE}$, $\mathcal{L}_\text{slope}$) dusuktir
cunku model izothermal icin zaten egitilmistir. Dolayisiyla optimizer termal
aktiviteyi bastirmak icin bir ``tesviki'' vardir --- az termal = az
termal loss = toplam loss dusmez ama momentum loss'lari minimize edilir.

Guard rail olarak $\mathcal{L}_{\theta,\text{min}}$ (Denklem~\ref{eq:loss-theta-min})
eklenmistir: $\theta_\text{rms} < 0.015$ durumunda penalize eder. Ancak bu
yapay bir esik, fiziksel termal yapiyi degil sadece minimum varyans
seviyesini zorlayarak ``hayatta kalmasi'' saglar. Model boylece
$\theta_\text{rms} \approx 0.015$'te sabitlenmekte ama \emph{fiziksel anlami
olan} termal yapi (plume'lar, termal sinir tabakasi) olusmamaktadir.


\subsubsection{$\langle v\theta \rangle$ Baglanti Cikmazi (Coupling Deadlock)}
\label{sec:vtheta-deadlock}

Termal flux $\langle v\theta \rangle$ hem hiz ($v$) hem sicaklik ($\theta$)
alaninin fonksiyonudur. $\mathcal{L}_{v\theta}$ loss'u iki yoldan
kucultulmeye calisilabilir:

\begin{enumerate}
  \item $\theta$'yi duzelt (termal yapi iyilestir) $\Rightarrow$
        \emph{dogru yol}
  \item $|v|$'yi artir (daha buyuk hiz = daha buyuk korelasyon) $\Rightarrow$
        \emph{yanlis yol}, TKE artar
\end{enumerate}

Optimizer tipik olarak 2.~yolu secmektedir cunku momentum parametrelerinin
gradient norm'u cok daha buyuktur ve $v$'yi artirmak daha kolay bir
optimizasyon yoludur. Sonuc: TKE referansin uzerine cikar
($\mathcal{L}_\text{TKE}$ artar), $\theta$ duzelmez, iki loss birbirine ters
calisir.

$\mathcal{L}_\text{corr}$ (Denklem~\ref{eq:loss-corr}) bu deadlock'u kirmak
icin tasarlanmistir: TKE'yi normalize ederek korelasyonun kendisini
hedefler. Ancak formulu genislettigimizde:

\begin{equation}
  C_{v\theta} = \frac{\langle v\theta \rangle}{\sqrt{\text{TKE}} \cdot \theta_\text{rms}}
  \label{eq:correlation-coefficient}
\end{equation}

\noindent $C_{v\theta}$'nin gradientleri \emph{yine} hem $v$ hem $\theta$'ya
akar. Quotient rule'dan:

\begin{equation}
  \frac{\partial C_{v\theta}}{\partial v}
  = \frac{\theta \cdot \sqrt{\text{TKE}} \cdot \theta_\text{rms}
          - \langle v\theta \rangle \cdot \frac{v}{2\sqrt{\text{TKE}}} \cdot \theta_\text{rms}}
         {(\sqrt{\text{TKE}} \cdot \theta_\text{rms})^2}
  \label{eq:corr-gradient}
\end{equation}

Bu ifade gostermektedir ki $C_{v\theta}$'nin gradient'i de momentum
parametrelerine akmaktan kacinamaz. Deadlock hafifletilir ama \emph{tamamen
cozulmez}.

\subsubsection{Buoyancy Instabilitesi}

\texttt{buoyancy\_strength} ogrenilebilir parametresi, gradient descent
tarafindan sistematik olarak arttirilmistir. Fiziksel mekanizma:

\begin{equation}
  \text{Ri}_\text{eff} \uparrow \implies
  \text{Ri}\,\theta\hat{\mathbf{e}}_y \uparrow \implies
  |v| \uparrow \implies
  \text{TKE} \uparrow
  \label{eq:buoyancy-feedback}
\end{equation}

TKE artisi $\mathcal{L}_\text{TKE,ref}$'i azalttigi icin optimizer bu
yonelimi ``faydali'' olarak degerlendirmektedir. Ancak fiziksel olarak
buoyancy'nin asiri artmasi sicaklik alanini bozar, tum termal yapi yikilir.
Clamp $[\text{0.001},\;\text{0.05}]$ uygulanmis olmakla birlikte ust sinira
yapisarak kalmistir.


% =============================================================================
\subsection{Basarisizligin Kok Neden Analizi}
\label{sec:root-cause}

Termal denemenin basarisizliginin on temel nedeni belirlenmistir. Bu
nedenler ``yapîsal'' (mimarinin temel siniri, kod degisikligi ile
cozulemez) ve ``muhendislik'' (prensipte duzeltilebilir, yeterli
muhendislik eforu ile) kategorilerine ayrilmaktadir.


\subsubsection{Yapisal Nedenler (Fundamental)}
\label{sec:fundamental-causes}

\paragraph{Neden 1: Skaler Parametrelerin Mekansal Yetersizligi}

INNATE'te her katmandaki termal difuzyon, $\kappa_\text{eff} =
\kappa + \nu_t/\text{Pr}_t$ formuluyle modellenmektedir. Burada
$\text{Pr}_t$ turbulansi Prandtl sayisidir ve her katmanda tek bir skaler
(\texttt{cs\_thermal}) ile parametrize edilmistir.

Gercekte $\text{Pr}_t$ \emph{mekansal olarak degisen} bir buyukluktur.
DNS verileri gostermektedir ki \cite{kays1994turbulent}:

\begin{itemize}
  \item Duvar yakininda: $\text{Pr}_t \to 1.1$--$1.3$
        (molekuler difuzyon baskisinda)
  \item Serbest kayma bolgesinde: $\text{Pr}_t \approx 0.7$--$0.9$
  \item Termal plume icerisinde: $\text{Pr}_t$ anizotropik ve $\sim 0.5$--$2.0$
        arasi degisir
\end{itemize}

Tek bir skaler ile bu mekansal heterojenitenin temsil edilmesi \emph{bilgi
teorik olarak imkansizdir}. INNATE katman basina tek bir
$\text{Pr}_t^\text{eff}$ atayabilir, bu da uzaysal ortalamadir:

\begin{equation}
  \text{Pr}_t^\text{layer} = \frac{1}{V}\int_V \text{Pr}_t(\mathbf{x})\,dV
  \label{eq:prt-average}
\end{equation}

Bu ortalama deger ne duvar yakinindaki ne de serbest akim bolgesindeki
dogru davranisi yakalar.

\emph{Cozum onerisi:} Adaptif mekansal $\text{Pr}_t$ --- ornegin spectral
space'de dalga sayisina bagimli $\text{Pr}_t(k)$, ya da fiziksel uzayda
yerel strain rate'e bagimli formulasyon.


\paragraph{Neden 2: $\langle v\theta \rangle$ Coupling Deadlock (Yapisal)}

Bolum~\ref{sec:vtheta-deadlock}'te aciklandigi uzere, $\langle v\theta \rangle$
buyuklugu hem momentum hem termal alana baglidir. 120 termal parametre ile
bu iki alanin \emph{birlesik} istatistiklerini kontrol etmek yapisal olarak
mumkun degildir: termal parametreler $\theta$'yi degistirir ama $v$'yi
dolayisiyla degistirir (buoyancy uzerinden), bu dolaylillik gradient sinyalini
zayiflatir.

Matematiksel olarak, Jacobian'in termal blogu:

\begin{equation}
  \frac{\partial \langle v\theta \rangle}{\partial \phi_\text{thermal}}
  = \left\langle v \frac{\partial \theta}{\partial \phi_\text{thermal}} \right\rangle
  + \left\langle \theta \frac{\partial v}{\partial \phi_\text{thermal}} \right\rangle
  \label{eq:vtheta-jacobian}
\end{equation}

\noindent Ikinci terim ($\partial v / \partial \phi_\text{thermal}$) buoyancy
zincirinden gecer: $\phi_\text{thermal} \to \theta \to \text{Ri}\theta\hat{e}_y
\to v$. Bu zincir 20 katman boyunca \emph{vanishing gradient} sorunu yasatir
cunku her katmanda diffusion ve projection operasyonlari gradient'i
sonumlendirir.


\paragraph{Neden 3: Multi-Scale Zaman Problemi}

Momentum ve termal alanlar farkli karakteristik zaman olceklerine sahiptir:

\begin{equation}
  \tau_\text{momentum} \sim \frac{L}{U} = \frac{L_y}{u_\text{rms}}
  \label{eq:tau-momentum}
\end{equation}

\begin{equation}
  \tau_\text{thermal} \sim \frac{L^2}{\kappa} = \text{Re} \cdot \text{Pr} \cdot L^2
  \label{eq:tau-thermal}
\end{equation}

$\text{Re} = 7{,}000$, $\text{Pr} = 0.71$ icin:
\begin{align*}
  \tau_m &\sim \frac{10}{0.13} \approx 77 \quad \text{(zaman birimi)} \\
  \tau_\theta &\sim 7000 \times 0.71 \times 100 = 4.97 \times 10^5
  \quad \text{(zaman birimi)}
\end{align*}

Oran: $\tau_\theta / \tau_m \sim 6{,}500$. INNATE'in 20 katmani
$20 \times \Delta t = 20 \times 0.025 = 0.5$ zaman birimi kapsamaktadir.
TBPTT (Truncated Backpropagation Through Time) ile 10--20 forward pass
yapilsa bile toplam pencere $\sim 10$ zaman birimi, bu momentum icin yeterli
ancak termal difuzyonun tam gelismesi icin \emph{dort mertebe yetersizdir}.

Bunun pratik sonucu: termal yapi henuz istatistiksel dengeye ulasmadan
loss hesaplaniyor ve gradient \emph{gecici} (transient) duruma gore
hesaplaniyor --- istatistiksel karararli (stationary) duruma degil.


\paragraph{Neden 4: Parametre Sayisi Yetersizligi}

120 termal parametre ile asagidaki uc fiziksel fenomenin \emph{hepsinin}
ayni anda ogrenilmesi gerekmektedir:

\begin{enumerate}
  \item \textbf{Buoyancy-turbulansi etkilesimi:} Richardson sayisina bagli,
        guclu nonlineerlik (turbulansi bastirilmasindan turbulansi
        guclendirilmesine gecis $\text{Ri} \sim 0.1$'de olur)
  \item \textbf{Termal sinir tabakasi:} $\text{Pr}_t$'nin mekansal degisimi,
        duvar-yakin davranis (periyodik BC'de gercek duvar yoktur ama
        buoyancy ``etkin duvarlar'' olusturur)
  \item \textbf{Mixed convection rejimleri:} Zorlanmis (forced) ve dogal
        (natural) konveksiyon arasindaki gecis
\end{enumerate}

Bu uc fizik birbirine nonlineer olarak baglidir. Ornegin buoyancy-turbulansi
etkilesimi, $\text{Ri}$ arttikca termal sinir tabakasini da degistirir,
bu da $\text{Pr}_t$ dagiliminini etkiler, bu da efektif $\kappa_t$'yi
degistirir ve tekrar buoyancy'yi etkiler. Bu geri besleme dongusu, skaler
parametrelerle ``duzenlenebilecek'' bir nonlineerligin otesindedir.

Karsilastirma: FNO \cite{li2021fno} ayni problem icin $\sim 500{,}000$
parametre kullanir ve bunlarin \emph{hepsi} mekansal bilgi tasir (Fourier
modlari). INNATE'in 120 termal parametresi bunun $\sim 1/4{,}000$'i kadardir.


\subsubsection{Muhendislik Nedenleri (Prensipte Duzeltilebilir)}
\label{sec:engineering-causes}

\paragraph{Neden 5: Damping Paradoksu}
\label{sec:damping-paradox}

INNATE'in referans verileri, \emph{sonumlu} (damped) LES cozucusunden
elde edilmistir. LES cozucusu buoyancy kaynakli instabiliteleri kontrol
etmek icin anizotropik sigma-profil sonumleme kullanmaktadir:

\begin{equation}
  \sigma(\mathbf{k}) = -\frac{1}{2}(\nu + \kappa)|\mathbf{k}|^2
  + \sqrt{\frac{1}{4}(\nu - \kappa)^2|\mathbf{k}|^4
  + \text{Ri}\frac{1}{L_y}\frac{k_x^2 + k_z^2}{|\mathbf{k}|^2}}
  \label{eq:growth-rate}
\end{equation}

\begin{equation}
  \gamma_\text{damp}(\mathbf{k}) = s \cdot \max(0, \sigma(\mathbf{k}))
  \label{eq:damping-profile}
\end{equation}

\noindent burada $s = 3.5$ guvenlik faktoru. Bu sonumleme momentum ($v$)
ve sicaklik ($\theta$) alanlarinin Fourier katsayilarina her adimda
uygulanmaktadir.

Sonuc olarak, LES referans Nusselt sayisi ``sonumlenmis'' bir fizigin
ciktisidir. Sonumlemesiz LES ile karsilastirildiginda:

\begin{table}[H]
  \centering
  \caption{Damping paradoksu: Re=10{,}000 LES karsilastirmasi.}
  \label{tab:damping-paradox}
  \begin{tabular}{lccc}
    \toprule
    & Nu & $\langle v\theta \rangle$ & $\theta_\text{rms}$ \\
    \midrule
    LES (sonumlu)       & 70.5   & $9.78 \times 10^{-4}$ & 0.0268 \\
    LES (sonumsuz)      & 124.2  & $1.74 \times 10^{-3}$ & 0.0550 \\
    Hata                 & \textbf{43.3\%} & \textbf{43.6\%} & \textbf{51.3\%} \\
    \bottomrule
  \end{tabular}
\end{table}

INNATE modeli sonumleme icermeyen ($s = 0$) bir PDE cozmektedir ama referans
olarak sonumlenmis degerleri hedeflemektedir. Bu sistematik \%43 uyumsuzluk,
$\mathcal{L}_\text{Nu,ref}$'in devre disi birakilmasina yol acmistir.
$\mathcal{L}_\text{Nu,ref}$ kapatilinca Nusselt sayisi \emph{hicbir loss
terimi tarafindan dogrudan yonlendirilmemektedir}; $\langle v\theta \rangle$
uzerinden dolayli hedeflenmesi ise yetersiz kalmistir.

\emph{Cozum onerisi:} Ya INNATE'e esit sonumleme eklemek ya da sonumlemesiz
LES referans uretmek. Ikincisi, buoyancy instabilitesi nedeniyle
$\text{Re} \geq 7{,}000$'de kararsizdir.


\paragraph{Neden 6: Gradient Routing Hatasi ($\text{Pr}_t$)}

Turbulansi Prandtl sayisi (\texttt{cs\_thermal}) EddyViscosity3D modulu
icerisinde tanimlanmistir. Gradient routing sistemi bu parametreyi
``momentum'' grubuna atamistir cunku EddyViscosity kelimesi momentum
keyword'lerinden birine eslesmektedir (\texttt{eddy}). Sonuc: termal loss'larin
gradient'i $\text{Pr}_t$'ye ulasmamistir.

Bu bug duzeltildikten sonra bile $\text{Pr}_t$ egitim surecinde $\sim 0.72$'de
takili kalmistir (beklenen: 0.85--1.0). Bunun nedeni Neden~1'dir:
tek bir skaler $\text{Pr}_t$, uzaysal dagilimi temsil edemedigi icin
optimizer tutarli bir gradient yonu bulamamaktadir.


\paragraph{Neden 7: LES Referans Verisi Kalitesi}

Periyodik sinir kosullarinda (periodic BC) ``gercek'' bir Nusselt sayisi
tanimlamak fiziksel olarak sorunludur. Nusselt sayisi normalde duvar isi
akisinin referans iletimine orani olarak tanimlanir:

\begin{equation}
  \text{Nu}_\text{wall} = \frac{q_\text{wall}}{k \cdot \Delta T / L}
  \label{eq:nu-wall}
\end{equation}

Periyodik kutuda duvar yoktur. Kullanilan tanimda (Denklem~\ref{eq:nusselt-def}),
$\text{Nu}$ hacimsel $\langle v\theta \rangle$'dan turetilmistir, bu da
sinir tabakasi fizigini dogrudan yansitmaz. LES referanslarindaki TKE ve
spektrum slope degerleri guvenilirdir (momentum buyuklukleri, sinir
kosulundan bagimsiz) ancak Nu, $\langle v\theta \rangle$ ve $\theta_\text{rms}$
degerleri sinir kosuluna duyarlidir.


\paragraph{Neden 8: Buoyancy Damping Uyumsuzlugu}

LES cozucusu \emph{Fourier uzayinda} her moda ozel sonumleme
(Denklem~\ref{eq:damping-profile}) uygulamaktadir. INNATE ise fiziksel
uzayda global bir sonumleme parametresi (\texttt{buoyancy\_damping\_strength})
kullanmaktadir:

\begin{equation}
  \text{INNATE:} \quad
  \begin{cases}
    v \leftarrow v - \gamma \cdot v \\
    \theta \leftarrow \theta - \gamma \cdot \theta
  \end{cases}
  , \quad \gamma = \text{BDS} \cdot \sqrt{\text{Ri}}
  \label{eq:innate-damping}
\end{equation}

Bu formul tum Fourier modlarina \emph{ayni} sonumlemeyi uygular. Oysa
fiziksel gerceklik, her dalga sayisinin farkli buyume hizina ($\sigma(\mathbf{k})$)
sahip olmasidir. Dusuk $|\mathbf{k}|$ modlari (buyuk olcek) yuksek
$|\mathbf{k}|$'ya (kucuk olcek) gore cok daha fazla sonumlenmelidir. Tek
parametreli global sonumleme ya buyuk olcekleri yeterince sonumleyemez (ve
model patlar) ya da kucuk olcekleri asiri sonumler (ve termal yapi kaybolur).


\paragraph{Neden 9: Elevator Mode ve $k_y = 0$ Filtreleme}

Periyodik kutuda $k_y = 0$ modlari (y-yonunde uniform, ``elevator'' modlari)
buoyancy etkisiyle lineer olarak buyur ve kontrolsuz enerji birikimi yaratir.
LES cozucu ve INNATE bu modlari her adimda sifirlamaktadir. Ancak bu
filtreleme sicaklik alaninin $k_y = 0$ bilesenini de tamamen yok eder.
Fiziksel olarak anlamsiz olmayan $k_y = 0$ termal modlar (ornegin yatay
homojen sicaklik degisimi) de budanmaktadir; bu durum $\theta_\text{rms}$
hesabini ve termal yapinin tam temsil edilmesini etkiler.


\paragraph{Neden 10: Organizasyonel Karar}

5-ajan teknik toplantisinda (ML Expert, CFD Expert, Architect, Codex Consultant,
Coordinator) asagidaki konsensus olusmustur:

\begin{quote}
  ``Termal 120 parametre ile bu mimaride CALISTIRILMAZ. Sorun bug degil,
  INNATE'in yapisal siniri. Skaler per-layer parametreler buoyancy coupling'in
  spatial heterojenigini yakalayamiyor. Termal birak, momentum parla, tez yaz.''
\end{quote}

Bu karar, kalan muhendislik eforu ile elde edilebilecek iyilesmenin
yapisal sinirlarin otesine gecilemeyecegi degerlendirmesine dayanmaktadir.


% =============================================================================
\subsection{Karsilastirmali Perspektif}
\label{sec:comparative}

INNATE'in termal basarisizligini daha genis bir baglamda
degerlendirmek icin, diger neural CFD yaklasimlariyla karsilastirma
sunulmaktadir.


\subsubsection{Fourier Neural Operator (FNO)}

FNO \cite{li2021fno} ve turevi olan FNO-3D, tum alanlari (hiz + sicaklik)
birlesik bir tensor olarak isler. Tipik parametre sayisi 500K--2M arasidir.
FNO'nun her katmani Fourier uzayinda \emph{mekansal bilgi tasir}: $R$ adet
Fourier modu icin ogrenilen $R \times R$ karmasik carpan matrisi, her noktadaki
etkilesimi implisit olarak yakalar. Bu yaklasim INNATE'ten temelden farklidir:
INNATE bir \emph{emulator} (ayni PDE'yi ogrenilebilir katsayilarla cozer),
FNO ise bir \emph{operator} (PDE'nin dogrudan cozum operatorunu ogrenir).

FNO termal akislar icin uygulandiginda \cite{li2024geometry}, sicaklik alani
momentum ile ayni $R$-modlu Fourier katmanindan gecer, boylece mekansal
coupling otomatik olarak ogrenilir. Ancak FNO'nun fizik bilgisi yoktur:
incompressibility, enerji korunumu gibi constraintler \emph{soft} loss olarak
empoze edilebilir ama garanti edilemez.


\subsubsection{DeepONet}

DeepONet \cite{lu2021deeponet} branch-trunk mimarisi ile parametrik PDE'leri
cozer. Multi-physics problemlerde tipik yaklasim, ayri branch network'ler
kullanmak ve trunk'ta paylasimli uzaysal temsil olusturmaktir. Termal
akislar icin rapor edilen parametre sayisi genellikle 100K--500K arasidir.

DeepONet'in termal coupling icin avantaji, branch network'lerin yeterli
kapasiteye sahip olmasidir: her parametre konfigurasyonu (Re, Ra, Pr) icin
farkli bir cozum ailesi ogrenilir. Dezavantaji, INNATE gibi fizik
garantisi tasimamasidir.


\subsubsection{MeshGraphNets}

Grafik tabanli yaklasimlar \cite{pfaff2021meshgraphnets} karmasik
geometrilerde basarili olmustur. Termal akislar icin node ozellikleri
$(u, v, w, p, \theta)$ seklinde genisletilir ve message passing ile hem
momentum hem termal etkilesimler ogrenilir. Tipik parametre sayisi 1--10M.

Bu modellerin INNATE'e gore temel avantaji, mekansal bilgiyi dogal olarak
tasimasadir: her node'un komsulariyla etkilesimi, yerel fizigi (termal sinir
tabakasi, plume yapisi) yakalayabilir.


\subsubsection{Karsilastirma Ozeti}

\begin{table}[H]
  \centering
  \caption{Neural CFD yaklasimlarinin termal akis basarimi karsilastirmasi.}
  \label{tab:comparison}
  \begin{tabular}{lcccl}
    \toprule
    Model & Param. & Fizik gar. & Termal & Yaklasim \\
    \midrule
    FNO-3D           & $\sim$500K  & Yok    & Basarili & Operator ogrenimi \\
    DeepONet         & $\sim$200K  & Yok    & Basarili & Branch-trunk \\
    MeshGraphNets    & $\sim$2M    & Yok    & Basarili & Message passing \\
    JAX-CFD          & $\sim$128K  & Kismi  & Basarili & Diferansiyel \\
    \textbf{INNATE}  & \textbf{287}& \textbf{Tam} & \textbf{Basarisiz} & Fizik noronlari \\
    \bottomrule
  \end{tabular}
\end{table}

Bu karsilastirma, INNATE'in termal basarisizliginin 287 parametreye ozgu
bir ``bedelinin'' oldugunu gostermektedir. Fizik garantileri (divergence-free,
enerji korunumu, CFL stabilitesi) karsiliginda mekansal temsil kapasitesi
feda edilmistir. Diger yaklasimlar yuz binlerce parametrenin mekansal
bilgi tasima kapasitesiyle termal coupling'i ogrenirken, INNATE'in skaler
parametreleri bu bilgiyi kodlayamamistir.

Bu bir tasarim trade-off'udur, basarisizlik degil: INNATE izothermal
turbulansi 287 parametre ile yakalayan \emph{ilk} modeldir. Termal
genisletme, bu ultra-kompakt mimarinin sinirini gostermistir.


% =============================================================================
\subsection{Alinan Dersler ve Gelecek Yonelimler}
\label{sec:lessons-future}

\subsubsection{Alinan Dersler}

\begin{enumerate}
  \item \textbf{Fizik garantisi $\neq$ fizik dogrulugu.}
        INNATE'in Leray projeksiyon ile \emph{garanti ettigi} divergence-free
        kosulu, momentum icin ``yeterli ve gerekli'' bir constrainttir.
        Ancak termal fizik icin boyle kesin constraintler yoktur; Nusselt
        sayisi, $\langle v\theta \rangle$ gibi buyuklukler \emph{istatistiksel}
        nicelliklerdir ve training sonucu ortaya cikar. Garanti
        mekanizmasinin yoklugunda, yeterli parametre kapasitesi sarttir.

  \item \textbf{Multi-objective loss dengelemesi, azinlik fizigini oldurur.}
        Gradient baskinligi sorunu (Neden 1), multi-task learning
        literaturunde \cite{chen2018gradnorm, yu2020gradient} bilinen bir
        problemdir. GradNorm, PCGrad gibi teknikler uygulanabilir; ancak
        INNATE'in 287-parametre boyutuyla bu tekniklerin overhead'i
        orantisizdir.

  \item \textbf{Referans veri ile model fiziginin tutarli olmasi sarttir.}
        Damping paradoksu (\%43 hata), verinin ``dogru'' olmasinin yetmedigini
        gostermektedir. Referans verisinin \emph{modelin cozdugu PDE ile
        ayni PDE'nin cozumu} olmasi gereklidir. Sonumlenmis LES $\neq$
        Sonumlenmemis INNATE.

  \item \textbf{Zaman olcegi uyumsuzlugu gizli bir katildir.}
        $\tau_\theta / \tau_m \sim 10^3$ orani, TBPTT penceresinin
        termal fizigin ogrenilebilmesi icin yetersiz kalmasina neden olur.
        Bu bulgu, gelecekteki multi-scale neural PDE tasarimlari icin
        genel bir uyaridir.

  \item \textbf{``Calismiyor'' da bilimsel bir sonuctur.}
        287-parametreli saf fizik noronlari ile mixed convection'in
        ogrenilememesi, modelin \emph{representational capacity} sinirini
        kesin olarak belirlemektedir. Bu bilgi, gelecek mimarilerin
        boyutlandirilmasi icin degerlidir.
\end{enumerate}

\subsubsection{Gelecek Yonelimler}

INNATE'in termal kapasitesini artirmak icin asagidaki yaklasimlar
onerilmektedir. Her biri, belirlenen yapisal nedenlerin en az birini
hedef almaktadir.

\begin{enumerate}
  \item \textbf{Adaptif Mekansal $\text{Pr}_t$:}
        Dalga sayisina bagimli $\text{Pr}_t(k) = a + b\,|k|^c$ formulasyonu
        ile mekansal heterojenitenin bir boyutunu yakalamak. Bu, Neden~1'i
        dogrudan hedefler. Ek parametre: $\sim 60$ (3 parametre $\times$ 20
        katman).

  \item \textbf{Buoyancy Sabitleme:}
        \texttt{buoyancy\_strength}'i ogrenilebilir parametre olmaktan
        cikarip $\text{Ri} = \text{Ra}/(\text{Re}^2\text{Pr})$ formulune
        sabitlemek. Bu, Neden~3 (buoyancy instabilitesi) ve kismen
        Neden~2'yi (deadlock) hafifletir.

  \item \textbf{Tutarli Referans Verisi:}
        INNATE'in \emph{birebir ayni} PDE'sini cozen (sonumlemesiz)
        bir DNS/LES cozucu ile referans uretmek. Bu, Neden~5'i
        (damping paradoksu) tamamen ortadan kaldirir. Ancak sonumlemesiz
        cozumun kararliligi $\text{Re} \geq 7{,}000$'de sorunludur.

  \item \textbf{LoRA Fine-Tuning:}
        Izothermal checkpoint'tan baslayarak, duzenli termal parametrelere
        dusuk rankli adaptasyon (Low-Rank Adaptation) eklemek. LoRA,
        izothermal bilgiyi korurken termal parametrelere ``ince ayar''
        kapasitesi kazandirir. $r = 4$ rank ile ek parametre: $\sim 320$
        (hala binler mertebesinde).

  \item \textbf{IMEX (Implicit-Explicit) Zaman Entegrasyonu:}
        Termal difuzyonu implicit, advection'i explicit yapmak.
        Bu, Neden~3 (multi-scale zaman) sorununu hafifletir cunku
        implicit adim CFL limitasyonu olmadan buyuk $\Delta t$ ile
        termal difuzyonu cozer.

  \item \textbf{Ikili Agirlik Egitimi (Two-Stage):}
        Faz~1: Momentum parametreleri ogrenilir, termal dondurulur (mevcut
        basarili izothermal rejim). Faz~2: Momentum dondurulur, \emph{sadece}
        termal parametreler ogrenilir. Bu, gradient baskinligini tamamen
        ortadan kaldirir (Neden~1 muhendislik boyutu).
        \texttt{thermal\_train.py} bu stratejiyi uygulamistir ancak
        yapisal nedenler (Nedenler~1--4) nedeniyle yetersiz kalmistir.
\end{enumerate}


\subsubsection{Neden Siniflandirmasi Ozet Tablosu}

\begin{table}[H]
  \centering
  \caption{Termal basarisizlik nedenlerinin siniflandirmasi.}
  \label{tab:cause-classification}
  \begin{tabular}{clcc}
    \toprule
    \# & Neden & Tip & Hedef cozum \\
    \midrule
    1 & Skaler $\text{Pr}_t$ yetersizligi & Yapisal & Adaptif $\text{Pr}_t(k)$ \\
    2 & $\langle v\theta \rangle$ deadlock & Yapisal & Ayrik termal network \\
    3 & Multi-scale zaman & Yapisal & IMEX entegrasyon \\
    4 & Parametre sayisi & Yapisal & LoRA / genisletme \\
    5 & Damping paradoksu & Muhendislik & Tutarli referans \\
    6 & Gradient routing hatasi & Muhendislik & Bug fix (\checkmark) \\
    7 & Referans verisi kalitesi & Muhendislik & Periodic BC analizi \\
    8 & Buoyancy damping uyumsuzlugu & Muhendislik & Fourier-domain damping \\
    9 & Elevator mode filtreleme & Muhendislik & Secici $k_y = 0$ \\
    10 & Organizasyonel karar & --- & Tez stratejisi \\
    \bottomrule
  \end{tabular}
\end{table}


\subsubsection{Sonuc}

INNATE'in termal denemesi, 287-parametreli saf fizik-noron mimarisinin
izothermal turbulansi basariyla emule edebildigini ancak mixed convection
icin yeterli temsil kapasitesine sahip olmadigini gostermistir. Bu sonuc,
modelin bir \emph{zayifligi} oldugu kadar bir \emph{karakterizasyonudur}:
literaturdeki ilk MLP-free neural CFD modeli olarak INNATE'in
calisabilirlik sinirini kesin olarak belirlenmistir.

Izothermal modda 167 etkin parametre ile elde edilen sonuclar ---
TKE $\sim$\%14 hata, spektrum slope $-1.7$, $\text{div}_\text{max} < 10^{-5}$
--- INNATE'in temel yaklasiminin gecerliligini kanitlamaktadir. Termal
genisletmenin basarisizligi, bu gecerliligin sinirini gostermekte ve
gelecek arastirmalar icin somut bir yol haritasi sunmaktadir.
